% ------------------------------------------------------------
% Page 92
% ------------------------------------------------------------

\textbf{En 1954}, Juliette Martin, mère de Gabriel Martin de Pourcarès reprend la propriété : elle avait
soigné avec dévouement Adeline, âgée et Marceau (estropié les deux propriétaires vivant sur
place. Elle espérait être aidée de son second mari Arthur Maurin et de son fils Gabriel né
d’un premier mariage. Or Arthur gardait un troupeau et était absent l’hiver. Quant à Gabriel,
sa femme préféra rester à Pourcarès près de l’école.

\textbf{1958-1979} \hspace{0.3cm} Elie Chanson, avec sa famille, s’installe à Connilergues. C’est en 1960 qu’un
WC avec fosse septique est aménagé..

\textbf{1979-2003} \hspace{0.3cm} Pierre Villaret, nouveau fermier se consacre à son troupeau de vaches et de
jeunes veaux (les broutards) abandonnant, peu ou prou la polyculture.
Sa femme est institutrice à l’école des Sœurs.

\textbf{2004} \hspace{1.25cm} Dominique Negron, nouveau fermier, soigne quarante cinq vaches et veaux. -
Sa femme travaille dans une boulangerie à Lasalle.

\vspace{0.4cm}

Poursuivant notre route, après de grands prés qui étaient tous à l’arrosage nous arrivons à
\textbf{Campis}.

Quelques maisons à l’arrivée d’un petit ruisseau dont celle de Pierre Bertrand, qui accusé de
guider et protéger pasteurs et prédicants fut condamné à vie, aux galères et mourut à l’hôpital
des forçats de Marseille à 29 ans, en 1698.

Un siècle plus tard fut assassiné, après avoir été torturé, le couple Poujol, dans la maison
encore appelée aujourd’hui la maison de « Poujol de Campis » par la troupe des brigands
royalistes recrutés par Jean Louis Solier dit : « Sans Peur » ancien prieur de Colognac près de
Lasalle.

Ensuite la vallée se rétrécit et le chemin se met à grimper. Nous laissons sur la droite
l’ancienne ferme en ruine de \textbf{Campredon}, au bord de la Brèze, qui fut, un temps, Maison
Forestière des Eaux et- Forêts, puis arrivons après trois ou quatre kilomètres aux \textbf{Oubrets}.

Il y a une trentaine d’années, grâce au Parc des Cévennes, la famille Vandermersch a repris la
totalité des terres. Un beau troupeau de chèvres, (complété par quelques brebis et plus
récemment par cinq vaches) donne une production régulière de fromages très appréciés.
Avec la location de quelques gîtes et bientôt l’accueil à la ferme la vie est possible aux
Oubrets. Les autres maisons du hameau ont été restaurées, en particulier l’ancien temple dont
nous verrons l’histoire ci-après ; en effet, deux récits de souvenirs nous sont parvenus : celui
d’un des derniers instituteurs des Oubrets, André Bazalgette, juste avant-t la dernière guerre,
en 1938 et, un peu plus anciennes les notes de Jean Gounelle qui fut pasteur de Meyrueis en
1933, en ces années d’avant-guerre où hameaux et villages résonnaient de cris d’enfants et de
sonnailles des troupeaux.
